\section{Conclusi\'on}

Este proyecto demostró que se pueden integrar datos separados de tráfico y 
siniestros usando un proceso ETL, y consolidarlos en un modelo de esquema de 
constelación. El principal valor analítico del trabajo está en la creación de 
métricas clave que superan la estadística básica, como el Índice de Eficiencia 
y el Índice de Riesgo. Estos índices permitieron evaluar el flujo de las 
avenidas y su peligro real, pero siempre de forma proporcional al volumen 
vehicular. Además, el análisis del Perfil Temporal mostró claramente los 
ciclos de movilidad a lo largo de la semana y el año.

Con estas bases, el trabajo logró ofrecer un enfoque orientado a la acción con 
la categorización de Puntos Críticos y el Score de Radar, herramientas 
diseñadas para priorizar objetivamente la necesidad de fiscalización en la 
ciudad. Aunque el procesamiento de datos mostró limitaciones en las fuentes 
oficiales (formatos inconsistentes y meses sin registros), los reportes 
interactivos finales transforman esta información en una base robusta para 
respaldar la toma de decisiones en planificación y seguridad vial.

